% from ega4-2-fr.tex
\oldpage[IV-2]{171}
II) \emph{Réduction au cas où $A$ est un anneau local intègre de dimension
1, $B$ un anneau de valuation discrète, clôture intégrale de $A$ et le
morphisme $\operatorname{Spec}(B)\to\operatorname{Spec}(A)$ radiciel.} ---
Soit $K$ le corps des fractions de $A$. On sait (0, 23.2.7) que $B$ est
intersection d'une famille $(V_\lambda)$ d'anneaux de valuation discrète
telle que, pour tout $x\in K$, on ait $x\in V_\lambda$ sauf pour un nombre
fini d'indices $\lambda$. On a donc une suite exacte de $A$-modules
\[
0\longrightarrow B\longrightarrow K\longrightarrow
\bigoplus_\lambda(K/V_\lambda).
\]
Posons $K'=K\otimes_A A'$, $V'_\lambda=V_\lambda\otimes_A A'$; par
platitude, on déduit de la suite exacte précédente une nouvelle suite
exacte
\[
0\longrightarrow B'\longrightarrow K'\longrightarrow
\bigoplus_\lambda(K'/V'_\lambda)
\]
et par suite $B'=\bigcap_\lambda V'_\lambda$. En outre
$\operatorname{Spec}(K')$ est la fibre du morphisme
$\operatorname{Spec}(A')\to\operatorname{Spec}(A)$ au point générique de
$\operatorname{Spec}(A)$, donc $K'$ est un anneau noethérien normal, et par
suite composé direct d'un nombre fini d'anneaux intègres et intégralement
clos; le composé direct $L'$ des corps des fractions de ces derniers est
l'anneau total des fractions de $A'$, donc aussi celui de $B'$. On peut donc
appliquer le lemme (6.14.1.1), et si l'on prouve que chacun des
$V'_\lambda$ est un anneau normal, il en sera de même de $B'$.

On sait d'autre part (0, 23.2.7) qu'on peut prendre les $V_\lambda$ tels
qu'il y ait une sous-$A$-algèbre finie $C$ de $K$ telle que $V_\lambda$
soit la clôture intégrale de $C_{\mathfrak p_\lambda}$, où
$\mathfrak p_\lambda$ est un idéal premier de hauteur 1 dans $C$. Si l'on
pose $C'_{\mathfrak p_\lambda}=C_{\mathfrak p_\lambda}\otimes_A A'$, le
morphisme $\operatorname{Spec}(C'_{\mathfrak p_\lambda})\to
\operatorname{Spec}(C_{\mathfrak p_\lambda})$ est normal (6.8.2) et l'on a
$V'_\lambda=V_\lambda\otimes_{C_{\mathfrak p_\lambda}}
C'_{\mathfrak p_\lambda}$; on peut donc remplacer $B$ par $V_\lambda$ et
$A$ par $C_{\mathfrak p_\lambda}$, donc supposer que $A$ est local, intègre
et de dimension 1, $B$ sa clôture intégrale et un anneau de valuation
discrète. Il y a une sous-$A$-algèbre finie $A_1$ de $B$ telle que le
morphisme $\operatorname{Spec}(B)\to\operatorname{Spec}(A_1)$ soit radiciel
(0, 23.2.5), ce qui entraîne en particulier que $A_1$ est aussi un anneau
local (évidemment de dimension 1); en outre on peut supposer que $B$ et
$A_1$ ont même corps résiduel (loc. cit.). On peut donc par la même méthode
remplacer $A$ par $A_1$. Appliquant si nécessaire le procédé du début de
I), on peut enfin supposer que $A'$ est lui aussi un anneau local et que
l'homomorphisme $A\to A'$ est local.

III) \emph{Fin de la démonstration.} --- Nous établirons d'abord le lemme
suivant:

\begin{lemma}[6.14.1.2]
\label{IV.2.s6.14.1.2-fr}
Soient $A$ un anneau local noethérien intègre de dimension 1, $A'$ un
anneau local noethérien, $A\to A'$ un homomorphisme local tel que le
morphisme correspondant $\operatorname{Spec}(A')\to\operatorname{Spec}(A)$
soit normal. Soient $K$ le corps des fractions de $A$, $\mathfrak m$
l'idéal maximal de $A$, $k=A/\mathfrak m$ son corps résiduel.
\begin{enumerate}
\item[(i)] Si $\mathfrak q'_j$ $(1\leqslant j\leqslant r)$ sont les idéaux
minimaux de $A'$, $K'=K\otimes_A A'$ est composé direct d'anneaux intègres
et intégralement clos $K'_j\supset A'/\mathfrak q'_j$
$(1\leqslant j\leqslant r)$, $K'_j$ ayant pour corps des fractions le corps
des fractions $L'_j$ de $A'/\mathfrak q'_j$, de sorte que l'anneau total des
fractions $L'$ de $A'$ s'identifie au composé direct des $L'_j$ et $A'$ à
un sous-anneau de $K'$.
\item[(ii)] L'idéal $\mathfrak p'=\mathfrak mA'$ de $A'$ est premier;
l'anneau $A'_{\mathfrak p'}$ est de dimension 1 et s'identifie à un
sous-anneau du produit des corps $L'_j$ pour les indices $j$ tels que
$\mathfrak q'_j\subset\mathfrak p'$.
\item[(iii)] Si $\overline{A'_{\mathfrak p'}}$ désigne le sous-anneau de
$L'$, produit de $A'_{\mathfrak p'}$ et des $L'_j$ tels que
$\mathfrak q'_j\not\subset\mathfrak p'$, on a
\begin{equation}
\tag{6.14.1.3}
\label{IV.2.s6.14.1.3-fr}
A'=K'\cap\overline{A'_{\mathfrak p'}}.
\end{equation}
\end{enumerate}
\end{lemma}

\begin{proof}
